\begin{definition}[The two native bonds in the charge-motion square]
    \label{def:peps_torus_charge_motion_bonds}
    \lean{TNLean.PEPS.torusChargeMotionWidthOne,
      TNLean.PEPS.torusChargeMotionHeightOne,
      TNLean.PEPS.torusPlaquetteChargeMotionBond,
      TNLean.PEPS.torusPlaquetteChargeMotionBond_val,
      TNLean.PEPS.torusPlaquetteChargeMotionBond_ne}
    \leanok
    \uses{def:peps_torus_plaquette_flux_geometry}
    On a finite torus with both periods at least three, the motion square
    based at $v=(x,y)$ has the four sites $(x,y),(x+1,y),(x+1,y+1),(x,y+1)$.
    Its distinct upper and lower charge bonds are the native right edges based at
    $(x,y+1)$ and $(x,y)$, respectively. Both ordered endpoints lie in
    the actual four-site region, including at either periodic seam.
    This is the diagram of \cite{Schuch2010PEPS},
    \texttt{eq:anyons:chargeon-move-setting}, lines 2489--2507.
\end{definition}

\begin{theorem}[Charge motion on the actual four original spins]
    \label{thm:peps_torus_physical_charge_motion}
    \lean{TNLean.PEPS.IsGIsometric.exists_unitary_torusPhysicalChargeMotion}
    \leanok
    \uses{def:peps_torus_charge_motion_bonds,
      thm:peps_regular_untwisted_physical_charge_motion}
    Let the original four-leg tensor be $G$-isometric for the regular action.
    At every translated position, one unitary on the four original spins
    moves the literal diagonal charge from the upper horizontal bond to
    the lower horizontal bond, retaining its character, internal parameter
    and all actual boundary labels. It is chosen before these data.
    All spanning-tree and root choices are derived internally.

    This is the finite-torus specialization of the charge-motion diagram in
    \cite{Schuch2010PEPS}, lines 2489--2507, with both periods at least three.
    The source proof further factors the operation into independent unitaries
    on the two columns; that factorization is not asserted here. No statement
    about energies or the parent Hamiltonian is made. The broader geometric
    scope remains separate in \cite{gap:rmp_peps_quantum_double_g_isometry}.
\end{theorem}
\begin{proof}\leanok
    The actual plaquette walk supplies connectivity of its induced four-site
    graph. The selected upper and lower edges belong to this region.
    Transport the given four-leg local isometry to its actual incident
    coordinates and apply the preceding original-spin motion theorem.
    At a seam the two endpoints are merely sorted differently: the identity
    bond retains the same shared charge label at both ends.
\end{proof}
